\documentclass[a4paper,11pt]{article}
\usepackage[utf8]{inputenc}
\usepackage{geometry}
\usepackage{xcolor}
\usepackage{listings}
\usepackage{hyperref}
\usepackage{graphicx}
\usepackage{amsmath}
\usepackage{amssymb}
\usepackage{tcolorbox}
\usepackage{float}
\usepackage{braket}
\usepackage{physics}

% Page geometry
\geometry{top=2.5cm, bottom=2.5cm, left=2.5cm, right=2.5cm}

% Custom colours
\definecolor{ibmblue}{RGB}{0, 45, 156}
\definecolor{ibmred}{RGB}{250, 75, 76}
\definecolor{purpleaccent}{RGB}{109, 40, 217}
\definecolor{amberaccent}{RGB}{245, 158, 11}
\definecolor{tealaccent}{RGB}{13, 148, 136}
\definecolor{codegreen}{rgb}{0,0.6,0}
\definecolor{codegray}{rgb}{0.5,0.5,0.5}
\definecolor{codepurple}{rgb}{0.58,0,0.82}
\definecolor{backcolour}{rgb}{0.95,0.95,0.92}

% Python code style
\lstdefinestyle{mystyle}{
    backgroundcolor=\color{backcolour},
    commentstyle=\color{codegreen},
    keywordstyle=\color{magenta},
    numberstyle=\tiny\color{codegray},
    stringstyle=\color{codepurple},
    basicstyle=\ttfamily\footnotesize,
    breakatwhitespace=false,
    breaklines=true,
    captionpos=b,
    keepspaces=true,
    numbers=left,
    numbersep=5pt,
    showspaces=false,
    showstringspaces=false,
    showtabs=false,
    tabsize=2
}
\lstset{style=mystyle}

% Custom boxes
\newtcolorbox{studentbox}{colback=blue!5!white,colframe=blue!75!black,title=Student Activity}
\newtcolorbox{theorybox}{colback=gray!5!white,colframe=gray!50!black,title=Theory Recap}
\newtcolorbox{warningbox}{colback=yellow!10!white,colframe=orange!80!black,title=Important}
\newtcolorbox{simulatorbox}{colback=purple!5!white,colframe=purple!75!black,title=Simulator}
\newtcolorbox{reflectionbox}{colback=teal!5!white,colframe=teal!75!black,title=Reflection Questions}
\newtcolorbox{connectionbox}{colback=green!5!white,colframe=green!60!black,title=Connections to Previous Labs}

% Bra-ket shorthands (fallback if braket/physics not available)
\providecommand{\ket}[1]{|#1\rangle}
\providecommand{\bra}[1]{\langle#1|}
\providecommand{\braket}[2]{\langle#1|#2\rangle}
\providecommand{\ketbra}[2]{|#1\rangle\langle#2|}

\title{\textbf{Laboratory Guide 5: Quantum Measurements}\\
\large Superposition, State Collapse, Spin \& Polarization on the Bloch Sphere}
\author{Course: Fundamental Concepts of Quantum Technologies}
\date{}

\begin{document}

\maketitle

\section{Introduction and Objectives}

This laboratory offers an \textbf{experimental verification and visualization} of the fundamental quantum
concepts presented in the \textit{Fundamental Concepts of Quantum Technologies} course. Unlike previous labs
that used Qiskit circuit execution, this laboratory uses \textbf{visual online simulators} to build physical
intuition for quantum measurement before introducing the algebraic formalism.

Students will explore superposition, quantum measurement, photon polarization, and electron spin as
represented on the Bloch sphere. Through guided simulations, they will analyze how measurement affects
quantum states, how probabilistic outcomes emerge, and how spin and polarization encode and manipulate
quantum information.

\vspace{0.3cm}
\noindent\textbf{Prerequisites:}
\begin{itemize}
    \item \textit{Fundamental Concepts of Quantum Technologies} (lecture series)
    \item Familiarity with Dirac notation: $\ket{\psi}$, $\bra{\phi}$, $\braket{\phi}{\psi}$
    \item Lab 1: \textit{Introduction to Qiskit} (optional, for companion notebook)
\end{itemize}

\begin{connectionbox}
\begin{itemize}
    \item \textbf{Lab 1 — Bloch Sphere:} In Lab 1 you watched the Bloch vector move when X, H, and Z gates were applied.
          Lab 5 revisits those same states, asking a different question: what happens when we \emph{observe} them?
    \item \textbf{Lab 1 — Exercise 3 (Phase \& Z gate):} The Z gate turned the Q-Sphere dot red but left
          the 50/50 probabilities unchanged. That phase was invisible because the measurement was in the Z-basis.
          Activity 2 here reveals why: the measurement basis determines which differences are detectable.
    \item \textbf{Lab 1 — Exercise 5 (Interference):} The H–Z–H sequence produced 100\% outcome $|1\rangle$
          by converting a phase into a measurable bit flip. Activity 5 shows that inserting a measurement
          between the H gates destroys that interference.
    \item \textbf{Lab 2 — Entanglement \& Measurement:} Measuring one qubit of a Bell state instantly
          determined the other. That correlation is the same state collapse studied here, applied to two qubits.
\end{itemize}
\end{connectionbox}

\noindent\textbf{Mode:} Online Simulators \quad\textbf{Credits:} 1 ECTS

\vspace{0.3cm}
\noindent\textbf{Learning Objectives:}
\begin{enumerate}
    \item Connect abstract quantum principles to measurable physical phenomena.
    \item Describe how measurement defines and changes information in quantum systems.
    \item Interpret visual representations of qubit states on the Bloch sphere.
    \item Explain how polarization and spin illustrate superposition and state collapse.
    \item Develop conceptual intuition about the probabilistic nature of quantum measurement.
\end{enumerate}

\newpage

\section{Theoretical Background}

\subsection{The Measurement Postulate}

Let $\ket{\psi}$ be a normalized qubit state and $\{|m\rangle\}_{m \in \{0,1\}}$ be an orthonormal
measurement basis. The \textbf{Born rule} gives the probability of obtaining outcome $m$:

\begin{equation}
    P(m) = |\braket{m}{\psi}|^2
    \label{eq:born}
\end{equation}

After the measurement, the state \textbf{collapses} to the corresponding eigenstate:

\begin{equation}
    |\psi\rangle \;\xrightarrow{\;\text{measure, outcome }m\;}\; |m\rangle
    \label{eq:collapse}
\end{equation}

This projection is \textbf{irreversible}: information about the pre-measurement state is generally lost.

\subsection{Bloch Sphere Parametrization}

Any single-qubit pure state can be written as:

\begin{equation}
    |\psi\rangle = \cos\!\left(\frac{\theta}{2}\right)|0\rangle
                 + e^{i\varphi}\sin\!\left(\frac{\theta}{2}\right)|1\rangle,
    \qquad \theta \in [0,\pi],\quad \varphi \in [0, 2\pi)
    \label{eq:bloch}
\end{equation}

The parameter $\theta$ is the polar angle from the north pole ($|0\rangle$) and $\varphi$ is the
azimuthal angle. The six cardinal states are:

\begin{center}
\begin{tabular}{lll}
\hline
\textbf{State} & \textbf{$(\theta,\varphi)$} & \textbf{Position} \\
\hline
$|0\rangle$ & $(0,\text{any})$ & North Pole \\
$|1\rangle$ & $(\pi,\text{any})$ & South Pole \\
$|{+}\rangle = \frac{|0\rangle+|1\rangle}{\sqrt{2}}$ & $(\pi/2, 0)$ & $+X$ axis \\
$|{-}\rangle = \frac{|0\rangle-|1\rangle}{\sqrt{2}}$ & $(\pi/2, \pi)$ & $-X$ axis \\
$|{i}\rangle = \frac{|0\rangle+i|1\rangle}{\sqrt{2}}$ & $(\pi/2, \pi/2)$ & $+Y$ axis \\
$|{-i}\rangle = \frac{|0\rangle-i|1\rangle}{\sqrt{2}}$ & $(\pi/2, 3\pi/2)$ & $-Y$ axis \\
\hline
\end{tabular}
\end{center}

\subsection{Pauli Operators and Their Eigenstates}

The three Pauli operators are the fundamental single-qubit observables:

\begin{equation}
    \hat{X} = \begin{pmatrix}0&1\\1&0\end{pmatrix},\qquad
    \hat{Y} = \begin{pmatrix}0&-i\\i&0\end{pmatrix},\qquad
    \hat{Z} = \begin{pmatrix}1&0\\0&-1\end{pmatrix}
    \label{eq:paulis}
\end{equation}

\noindent Measuring in the $Z$-basis collapses the state to $|0\rangle$ or $|1\rangle$ (north/south poles).\\
Measuring in the $X$-basis collapses the state to $|{+}\rangle$ or $|{-}\rangle$ (equatorial $\pm X$).\\
Measuring in the $Y$-basis collapses the state to $|{i}\rangle$ or $|{-i}\rangle$ (equatorial $\pm Y$).

\subsection{Non-Commutativity of Pauli Observables}

\begin{theorybox}
The commutator $[\hat{X}, \hat{Z}]$ quantifies how much the order of operations matters:
\begin{align}
    [\hat{X}, \hat{Z}] &= \hat{X}\hat{Z} - \hat{Z}\hat{X}
    = \begin{pmatrix}0&1\\1&0\end{pmatrix}\begin{pmatrix}1&0\\0&-1\end{pmatrix}
    - \begin{pmatrix}1&0\\0&-1\end{pmatrix}\begin{pmatrix}0&1\\1&0\end{pmatrix} \nonumber\\
    &= \begin{pmatrix}0&-1\\1&0\end{pmatrix} - \begin{pmatrix}0&1\\-1&0\end{pmatrix}
    = \begin{pmatrix}0&-2\\2&0\end{pmatrix} \neq 0
    \label{eq:commutator}
\end{align}
Since $[\hat{X}, \hat{Z}] \neq 0$, the observables $\hat{X}$ and $\hat{Z}$ are \textbf{incompatible}: they
share no common eigenstates and cannot be simultaneously measured without disturbing each other.
This is the algebraic root of the sequential measurement randomness observed in Activity 5.
\end{theorybox}

\subsection{Photon Polarization as a Qubit}

A photon's polarization degree of freedom maps directly to a qubit:
\begin{align}
    \text{Horizontal } H &\;\leftrightarrow\; |0\rangle, \\
    \text{Vertical } V &\;\leftrightarrow\; |1\rangle, \\
    \text{Diagonal } (+45°) &\;\leftrightarrow\; |{+}\rangle = \tfrac{1}{\sqrt{2}}(|0\rangle+|1\rangle).
\end{align}
A polarizer at angle $\alpha$ performs a projective measurement in the basis
$\{|\alpha\rangle, |\alpha+90°\rangle\}$. The transmission probability for a photon
in polarization state $|\theta_0\rangle$ is:
\begin{equation}
    P(\text{pass}) = |\langle\alpha|\theta_0\rangle|^2 = \cos^2(\alpha - \theta_0)
    \qquad \text{(Malus's Law / Born rule)}
    \label{eq:malus}
\end{equation}

\subsection{Electron Spin and the Stern-Gerlach Experiment}

An electron's spin-$\frac{1}{2}$ is a two-level quantum system — identical in structure to a qubit.
The correspondence with the Bloch sphere is:
\begin{center}
\begin{tabular}{lll}
\hline
\textbf{Spin state} & \textbf{Qubit state} & \textbf{Bloch position}\\
\hline
Spin-up along $Z$ & $|0\rangle$ & North Pole \\
Spin-down along $Z$ & $|1\rangle$ & South Pole \\
Spin-up along $X$ & $|{+}\rangle$ & $+X$ equator \\
Spin-up along $Y$ & $|i\rangle$ & $+Y$ equator \\
\hline
\end{tabular}
\end{center}
The Stern-Gerlach apparatus oriented along a given axis measures the spin component along that axis,
splitting the beam into two outputs (spin-up and spin-down). The probability of each output follows
the Born rule: $P(\uparrow_{\hat{n}}) = |\langle\uparrow_{\hat{n}}|\psi\rangle|^2$.
This provides a direct physical realization of quantum measurement and state collapse.

\newpage

\section{Simulator Tools}

\begin{simulatorbox}
The following external simulators are used. Open each in a browser tab — no installation required.
\begin{itemize}
    \item \textbf{QuVis Bloch Sphere} (Activities 1--2): \\
          \url{https://www.st-andrews.ac.uk/physics/quvis/}
    \item \textbf{Quantum Flytrap Virtual Lab} (Activity 3): \\
          \url{https://lab.quantumflytrap.com/}
    \item \textbf{PhET Stern-Gerlach} (Activity 4): \\
          \url{https://phet.colorado.edu/en/simulations/stern-gerlach}
    \item \textbf{IBM Quantum Composer} (Activity 5): \\
          \url{https://quantum.ibm.com/composer}
\end{itemize}
\end{simulatorbox}

\newpage

\section{Activities}

\subsection{Activity 1: State Preparation \& Bloch Sphere}

\textbf{Learning goal:} Build visual intuition for qubit states by placing them on the Bloch sphere
and understanding the spherical coordinate parametrization $(\theta,\varphi)$.

\textit{Connection to Lab 1: In Lab 1, X, H, and Z gates moved the Bloch vector between specific positions.
Here you will explore those positions directly, without any gate, understanding the coordinates $(\theta,\varphi)$
of the states you already encountered.}

\begin{studentbox}
\textbf{Tool:} QuVis Bloch Sphere Simulator.

\begin{enumerate}
    \item Prepare each of the six cardinal states
          ($|0\rangle$, $|1\rangle$, $|{+}\rangle$, $|{-}\rangle$, $|i\rangle$, $|{-i}\rangle$)
          by adjusting the $\theta$ and $\varphi$ sliders.
    \item Observe how the Bloch vector moves to the poles for $|0\rangle$ and $|1\rangle$,
          and to the equator for the remaining four states.
    \item Attempt to prepare the state
          $|\psi\rangle = \cos(30°)|0\rangle + e^{i\pi/3}\sin(30°)|1\rangle$ ($\theta = 60°$, $\varphi = 60°$).
    \item Note that rotating $\varphi$ while keeping $\theta = \pi/2$ traces a circle on the equator.
          This corresponds to changing the relative phase without changing measurement probabilities in the $Z$-basis.
\end{enumerate}

\textbf{Expected observation:} Global phase is not visible on the Bloch sphere.
The state $e^{i\delta}|\psi\rangle$ looks identical to $|\psi\rangle$.
\end{studentbox}

\begin{reflectionbox}
\textbf{Question 1:} Why do $|{+}\rangle$ and $|{-}\rangle$ both lie on the equator ($\theta = \pi/2$)
but on opposite sides ($\varphi = 0$ vs $\varphi = \pi$)?
What does $\varphi$ physically represent?
\end{reflectionbox}

\subsection{Activity 2: Measurement in Different Bases}

\textbf{Learning goal:} Measure the same state in the $Z$, $X$, and $Y$ bases and observe how the
measurement basis determines the outcome probability distribution.

\textit{Connection to Lab 1 — Exercise 3: Applying the Z gate after H left the 50/50 probabilities
unchanged, even though the Q-Sphere dot turned red. That phase change was invisible because the
measurement was in the Z-basis, which cannot distinguish $|{+}\rangle$ from $|{-}\rangle$ by probabilities.
Measuring in the X-basis makes that difference fully visible.}

\begin{studentbox}
\textbf{Tool:} QuVis Bloch Sphere Simulator.

\begin{enumerate}
    \item Prepare the state $|{+}\rangle$ ($\theta = \pi/2$, $\varphi = 0$).
    \item Select the \textbf{Z-basis} and click ``Measure'' 10 times. Record the outcomes.
          Verify: approximately $P(0) = P(1) = 50\%$.
    \item Reset the state to $|{+}\rangle$. Select the \textbf{X-basis} and measure.
          Verify: $P({+}) = 100\%$ every time. \textit{Why?} Because
          $|\langle{+}|{+}\rangle|^2 = 1$.
    \item Reset to $|{+}\rangle$. Select the \textbf{Y-basis}.
          Verify: $P({i}) = P({-i}) = 50\%$.
    \item Repeat the entire sequence for the state $|{-}\rangle$.
\end{enumerate}

\textbf{For each trial, record:}
\[
\text{Initial state} \;\longrightarrow\; \text{Basis} \;\longrightarrow\; P(0), P(1) \;\longrightarrow\; \text{Observed outcome}
\]
\end{studentbox}

\begin{reflectionbox}
\textbf{Question 2:} Why does measuring $|{+}\rangle$ in the X-basis always give the same result?
Compute $|\langle{+}|{+}\rangle|^2$ and $|\langle{-}|{+}\rangle|^2$ explicitly using equation~\eqref{eq:born}.
\end{reflectionbox}

\subsection{Activity 3: Photon Polarization}

\textbf{Learning goal:} Map classical photon polarization to qubit states. Observe how a polarizer acts
as a quantum measurement device and causes state collapse.

\begin{studentbox}
\textbf{Tool:} Quantum Flytrap Virtual Lab.

\begin{enumerate}
    \item Prepare a \textbf{horizontally polarized} photon ($|0\rangle$) and pass it through a
          \textbf{vertical} polarizer ($|1\rangle$ basis). Observe: 0\% transmission.
    \item Prepare a \textbf{diagonally polarized} photon ($|{+}\rangle$) and measure with a
          vertical polarizer. Observe: 50\% transmission (each photon either passes or does not).
    \item Take a photon that \textit{passed} through the vertical polarizer in step~2 (now in state $|1\rangle$).
          Measure again with a \textbf{horizontal} polarizer. Observe: 0\% — the state collapsed to $|1\rangle$.
    \item Insert a $45°$ polarizer \textit{between} two crossed ($H$/$V$) polarizers.
          Observe that some light now passes through the final $V$ polarizer.
          Explain why using state collapse.
\end{enumerate}
\end{studentbox}

\begin{reflectionbox}
\textbf{Question 3:} A diagonally polarized photon hits a vertical polarizer and passes through.
What is its polarization state immediately after the polarizer?
Use equation~\eqref{eq:collapse} to justify your answer.
\\[4pt]
\textbf{Question 4:} Using equation~\eqref{eq:malus} (Malus's Law / Born rule), compute the
transmission probability for a $45°$ photon through a $30°$ polarizer.
\end{reflectionbox}

\subsection{Activity 4: Electron Spin \& Stern-Gerlach}

\textbf{Learning goal:} Use the Stern-Gerlach apparatus to observe spin quantization and understand
how spin-up and spin-down states along different axes map to Bloch sphere positions.

\begin{studentbox}
\textbf{Tool:} PhET Stern-Gerlach Simulation.

\begin{enumerate}
    \item Send atoms through a \textbf{Z-oriented} apparatus. Observe two output beams.
    \item Select only the spin-up output ($|\!\uparrow_z\rangle = |0\rangle$) and feed it
          into a \textbf{second Z apparatus}. Observe: 100\% spin-up every time.
          \textit{The state is now definite in $Z$.}
    \item Instead, feed the spin-up beam into an \textbf{X-oriented} apparatus.
          Observe: 50/50 split. Why? Expand:
          \[
          |0\rangle = \frac{|{+}\rangle + |{-}\rangle}{\sqrt{2}}
          \]
    \item Select only $|{+}\rangle$ from the X apparatus. Feed into a \textbf{third Z apparatus}.
          Observe: 50/50 again — the X measurement erased the Z information.
    \item Document each step with a sketch of the apparatus and the outcome fractions.
\end{enumerate}
\end{studentbox}

\begin{reflectionbox}
\textbf{Question 5:} Use the commutator result from equation~\eqref{eq:commutator} to explain
algebraically why the X measurement in step~3 destroys the Z information established in step~2.

\textbf{Question 6:} What would you observe in step~4 if $[\hat{S}_x, \hat{S}_z] = 0$?
\end{reflectionbox}

\subsection{Activity 5: Sequential Measurements \& State Collapse}

\textbf{Learning goal:} Demonstrate that the sequence $Z \to X \to Z$ produces random outcomes in
the final Z measurement, even when the first Z gave a definite result. This concretizes $[\hat{X},\hat{Z}] \neq 0$.

\textit{Connection to Lab 1 — Exercise 5: The H–Z–H sequence gave 100\% outcome $|1\rangle$ because
the two H gates converted the invisible phase into a measurable bit flip. Inserting a measurement
between the H gates (as in this activity) breaks that interference and reintroduces randomness.}

\begin{studentbox}
\textbf{Tool:} IBM Quantum Composer.

\begin{enumerate}
    \item Start with $|0\rangle$. Add a measurement gate. Run 1000 shots.
          Verify: 100\% outcome $0$.
    \item Reset. Add $H$, then measure ($X$-basis equivalent). Observe: 50/50.
    \item Reset. Start with $|0\rangle$, add $H$ then measure (X), then $H$ again then measure (Z).
          Observe: the final Z measurement gives 50/50.
          Interpretation: measuring in X collapses the state to $|{+}\rangle$ or $|{-}\rangle$;
          applying $H$ converts these to $|0\rangle$ or $|1\rangle$ — each with probability $\frac{1}{2}$.
    \item \textbf{Control experiment:} $Z \to Z \to Z$. Start from $|0\rangle$, measure three times in $Z$.
          Observe: always $0$. The Z basis is stable under Z measurements.
\end{enumerate}

\begin{warningbox}
\textbf{Qiskit / IBM Composer Convention (Little Endian):}
In a bitstring like $|10\rangle$, the \textbf{rightmost} bit represents qubit $q_0$ and the \textbf{leftmost}
bit represents qubit $q_1$. Keep this in mind when reading multi-qubit measurement results.
\end{warningbox}
\end{studentbox}

\begin{reflectionbox}
\textbf{Question 7:} After the X measurement, the state has collapsed to either $|{+}\rangle$ or $|{-}\rangle$.
Both are equal superpositions of $|0\rangle$ and $|1\rangle$.
Can the subsequent Z measurement ever recover the original Z information? Why not?

\textbf{Question 8:} Compare Z$\to$Z$\to$Z with Z$\to$X$\to$Z.
Why does the first sequence give a deterministic final Z result, while the second does not?
\end{reflectionbox}

\newpage

\section{Assessment}

\textbf{Question A (Born Rule):}
A qubit is prepared in state $|\psi\rangle = \frac{\sqrt{3}}{2}|0\rangle + \frac{1}{2}|1\rangle$.
Compute:
\begin{enumerate}
    \item $P(|0\rangle)$ and $P(|1\rangle)$ for a Z-basis measurement.
    \item $P(|{+}\rangle)$ and $P(|{-}\rangle)$ for an X-basis measurement.
    \item Verify that probabilities sum to 1 in each case.
\end{enumerate}

\textbf{Question B (State Collapse):}
After measuring $|\psi\rangle = |{+}\rangle$ in the Z-basis and obtaining outcome $0$:
\begin{enumerate}
    \item What is the post-measurement state?
    \item If the same qubit is immediately measured again in the Z-basis, what is $P(0)$?
    \item If instead it is measured in the X-basis, what is $P({+})$?
\end{enumerate}

\textbf{Question C (Non-Commutativity):}
You prepare $|0\rangle$, measure in Z (outcome: 0), measure in X (outcome: ${+}$), then measure in Z.
\begin{enumerate}
    \item What is the state after each measurement?
    \item What are the final Z probabilities? Show with the Born rule.
    \item Explain why Z$\to$Z$\to$Z gives a deterministic result but Z$\to$X$\to$Z does not.
\end{enumerate}

\section{Resources}

\begin{itemize}
    \item Nielsen \& Chuang, \textit{Quantum Computation and Quantum Information}, Ch.~2.2 (Measurement Postulate).
    \item Griffiths, \textit{Introduction to Quantum Mechanics}, Ch.~1 (Statistical Interpretation).
    \item Preskill, \textit{Lecture Notes on Quantum Information}, Ch.~2 (Foundations).
    \item Qiskit Documentation: \url{https://qiskit.org/documentation/} — Statevector \& measurement simulation.
    \item Companion Jupyter notebook: \texttt{quantum-labs-2026/lab5/lab5\_quantum\_measurements.ipynb}.
\end{itemize}

\end{document}
